\begin{figure}[tb]
\centering
\begin{tikzpicture}[
  box/.style={draw=blue!45!black,rounded corners=2pt,fill=blue!3,
              align=center,text width=5.2cm,minimum height=1.3cm,inner sep=6pt},
  >={Stealth[length=2mm]},font=\small]
  \node[box] (independent) at (0,0)
    {Unconditioned hole picture $\mathbb Q_{\mathbf a}$\\
     $K=\mathbb E_{\mathbb Q_{\mathbf a}}\Pi$};
  \node[box] (conditioned) at (7,0)
    {Conditioned hole picture\\
     $\mathbb Q_{\mathbf a}(\,\cdot\mid G_J)$};
  \node[box] (walk) at (0,-2.6)
    {Markov walk with transition $K$\\
     First positive $J$-hitting kernel};
  \node[box] (moment) at (7,-2.6)
    {Conditional next-pin mean\\
     $\mathbb E[\Sigma\mid G_J]$};
  \draw[->,thick] (independent) -- node[above,align=center]
    {condition\\on $G_J$} (conditioned);
  \draw[->,thick] (independent) -- node[left,align=right]
    {retain the\\one-step mean} (walk);
  \draw[->,thick] (conditioned) -- node[right,align=left]
    {retain the\\next-pin mean} (moment);
  \draw[<->,thick] (walk) -- node[below,align=center]
    {equal\\kernels} (moment);
\end{tikzpicture}
\caption{The first-moment reduction at a fixed exposure and a finite set $J$
of pins on selected time cuts, with $\mathbb Q_{\mathbf a}(G_J)>0$.  The exterior-moment identity is used for the unconditioned
law in the upper-left box; Proposition~\ref{prop:first-moment} proves
that the two routes to the lower-right box agree.  The lower-left walk
resamples its transitions, whereas the upper-right law follows a single
sampled permutation.  The equality concerns the next-pin mean and does
not identify these two path laws.  When the exposure is averaged, its
actual marginal still contains the factor $g_P(\mathbf a)$ in
\eqref{eq:exposure-outer-law}.}
\label{fig:exposure-first-moment}
\end{figure}
